% Imporditud: tools/import_eksperimendid.py
% Allikas: Eesti füüsikaolümpiaadi eksperimendi ülesannete
%   kogu (Rael Kalda, 2023), ülesanne Ü10, lahendus L10.
\setAuthor{EFO žürii}
\setRound{piirkonnavoor}
\setYear{2003}
\setNumber{P E1}
\setDifficulty{1}
\setTopic{vedelike mehaanika}
\setVahendid{väike silindriline plastmasstopsik, anum veega, kaks 10-sendist ja üks 1-kroonine, joonlaud, permanentne viltpliiats}
\setPunktid{10}
\setAllikas{Eksperimendi kogu Ü10, lahendus lk 37}
\setLahendusseis{toores}

\prob{Ühekroonine}
Leida ühekroonise metallraha mass.

\hint

\solu
Võtame topsikusse veidike vett ja asetame sinna põhja 10-sendised nii, et topsik ujub stabiilselt püstises asendis. Märgime viltpliiatsiga topsi seinale ujumissügavuse. Lisame topsi ühekroonise ja märgime uue ujumissügavuse. Mõõdame kõrguste vahe h ja topsi välisdiameetri d. Mündi mass m = $\pi \cdot$ h $\cdot$ $d_2$ $\cdot \varrho$/4, kus $\varrho$ on vee tihedus.

Hindamine: Idee (3 p), tehniline teostus ja mõõtmine (3 p), mündi massi arvutamise valem (3 p), arvuline tulemus m = 5 $\pm$ 1,5 g (1 p).
\probend
